% SafeWard, AI4015 Phase 2 paper (IEEEtran, two-column, 8-10 pages incl. references).
% Upload this folder to Overleaf (New Project -> Upload Project -> zip of paper/).
%
% STATUS
% - Sections 1-7 contain a DRAFT prepared with AI assistance under the instructor's written
%   permission. The group must read, correct and revise every paragraph, and check every claim
%   against the repository before submission.
% - \tbd{...} marks numbers or facts that must be filled from results/ once experiments finish.
% - Sections 8-12 still hold CHECKLIST comments; they are drafted after results exist.
% - Every number must come from results/ (brief Sec. 5). Tables marked AUTO are \input from
%   results/tables/ and must not be retyped.

\documentclass[conference]{IEEEtran}
\usepackage{cite}
\usepackage{amsmath,amssymb}
\usepackage{graphicx}
\usepackage{booktabs}
\usepackage{multirow}
\usepackage{xcolor}
\usepackage{url}
\usepackage[hidelinks]{hyperref}

\newcommand{\tbd}[1]{\textcolor{red}{[TBD: #1]}}

\begin{document}

\title{SafeWard: A Triage-Aware, Safety-Guarded Multi-Agent Hospital Simulation for Clinical Reasoning}

\author{
\IEEEauthorblockN{Hudiya e Zehra}
\IEEEauthorblockA{23I-0761\\ \textit{i230761@isb.nu.edu.pk}}
\and
\IEEEauthorblockN{Aleeza Azhar}
\IEEEauthorblockA{23I-0765\\ \textit{i230765@isb.nu.edu.pk}}
\and
\IEEEauthorblockN{Abdul Ghufran}
\IEEEauthorblockA{23I-0576\\ \textit{i230576@isb.nu.edu.pk}}
\\
\IEEEauthorblockA{Project G2-04, AI4015 Agentic Artificial Intelligence, FAST-NUCES Islamabad}
}

\maketitle

% =====================================================================
\begin{abstract}
% Drafted last, once the gap-resolution table is final (150-250 words): problem, approach,
% data, one headline number per gap, contribution.
\tbd{abstract}
\end{abstract}

\begin{IEEEkeywords}
multi-agent systems, large language models, clinical decision support, emergency triage, AI safety, agent memory
\end{IEEEkeywords}

% =====================================================================
\section{Introduction}

Large language models (LLMs) now get to expert-level scores on static medical question-answering
benchmarks such as MedQA~\cite{jin2021medqa,singhal2023medpalm,nori2023medprompt}. In these
benchmarks the model get the complete case at once. Clinical work is unlike: information
is incomplete at first and is acquired through questions, examinations and tests. Interactive
simulators address this mismatch by letting a doctor agent consult a simulated patient over
several turns~\cite{schmidgall2024agentclinic,fan2025aihospital,tu2025amie}. MedAgentSim~\cite{almansoori2025medagentsim},
our base paper, extends this line with a measurement agent that returns test results only on
request and with a self-improving memory of earlier cases. With LLaMA~3.3~70B it reports a
diagnostic accuracy on MedQA that rises from 54.7\% for its multi-agent baseline to 70.8\% for
the full system.

These simulators share an implicit assumption: every patient walks directly into a one-to-one
consultation, and the only outcome that matters is the final diagnosis. Emergency care starts
earlier and fails in other ways. Patients are first triaged by urgency and queued, and two of
the most harmful errors occur outside the diagnosis itself. \emph{Under-triage} assigns a
critically ill patient a less urgent level, delaying care; recent studies report that LLMs tend
to under-triage high-acuity cases~\cite{npjai2026emergency,lee2025edtriagebias}. \emph{Unsafe actions},
such as a contraindicated drug, a missed red flag, or compliance with a patient's false premise,
cause harm even when the diagnosis is right; LLMs comply with illogical medical requests at
rates of up to 100\%~\cite{chen2025sycophancy}. Finally, self-evolving memory introduces a risk
of its own: a stored lesson that is wrong can be retrieved for later, similar patients and
spread the error, a weakness that has been exploited deliberately in poisoning attacks on agent
memory~\cite{chen2024agentpoison}.

We therefore study three gaps in MedAgentSim. \textbf{G1 (Triage):} patients are not triaged or
queued. \textbf{G2 (Safety):} unsafe orders and sycophantic compliance are never measured.
\textbf{G3 (Verified memory):} lessons enter memory without verification. We ask three research
questions:
\begin{itemize}
  \item \textbf{RQ1:} Does an interviewing Triage Nurse agent reduce under-triage of high-acuity
  patients compared with a single LLM given the full record and with rule-based triage?
  \item \textbf{RQ2:} Does an independent Safety Officer reduce unsafe actions more than doctor
  self-review, without degrading diagnostic accuracy?
  \item \textbf{RQ3:} Does gating memory writes on verified outcomes prevent injected wrong lessons
  from spreading to similar cases?
\end{itemize}

Our contributions are as follows:
\begin{enumerate}
  \item \textbf{SafeWard}, an open-source multi-agent hospital simulation that keeps the MedAgentSim
  consultation loop and adds triage, an acuity-ordered queue, an independent safety review and a
  verified memory, each of which can be switched off for ablation.
  \item An evaluation protocol that reports triage, safety and memory-integrity metrics next to
  diagnostic accuracy, using the KTAS emergency triage dataset~\cite{moon2019ktas}, an adversarial
  safety set and a memory-poisoning set.
  \item An empirical study on fixed splits, with repeated runs where the budget allows, that tests
  one pre-registered hypothesis per gap and reports positive, partial and negative results.
\end{enumerate}

% =====================================================================
\section{Related Work}

\textbf{Medical LLMs and static benchmarks.} Domain-tuned and prompted foundation models achieve
high accuracy on licensing-exam style questions~\cite{singhal2023medpalm,nori2023medprompt}, with
MedQA~\cite{jin2021medqa} as the most common benchmark. These evaluations present the full case
at once and score a single answer, so they do not test whether a model can gather information.

\textbf{Interactive clinical simulation.} AgentClinic~\cite{schmidgall2024agentclinic} converts MedQA,
NEJM and MIMIC-IV~\cite{johnson2023mimiciv} cases into objective structured clinical examinations in
which a doctor agent must question a patient agent and request tests from a measurement agent.
AI Hospital~\cite{fan2025aihospital} evaluates multi-turn consultations against scripted
non-player characters and shows that performance drops sharply compared with the full-case
setting. AMIE~\cite{tu2025amie} trains a diagnostic dialogue model by self-play and outperforms
primary-care physicians in a randomised text-consultation study, but the model and data are
closed. Agent Hospital~\cite{li2024agenthospital} and MedAgentSim~\cite{almansoori2025medagentsim}
let agents improve across patients; MedAgentSim does so with kNN retrieval over past cases,
reflection on errors, chain-of-thought prompting~\cite{wei2022cot} and majority voting. Building on
generative agents~\cite{park2023generative}, these systems model the consultation well, yet none
models triage, a patient queue or safety review, and all score diagnostic accuracy only.

\textbf{Multi-agent medical reasoning.} MedAgents~\cite{tang2024medagents} and MAC~\cite{chen2025mac}
show that discussion among several doctor agents can improve diagnosis, and MMedAgent-RL~\cite{xia2026mmedagentrl}
learns a triage-and-referral collaboration with reinforcement learning. MedAgentBoard~\cite{zhu2025medagentboard}
cautions that multi-agent systems do not reliably beat a strong single LLM and that their extra
cost must be justified. We follow this advice by comparing each agent against single-model and
rule-based baselines on the same cases.

\textbf{LLMs for emergency triage.} On the Korean Triage and Acuity Scale (KTAS)~\cite{moon2019ktas}
and the Emergency Severity Index~\cite{gilboy2020esi}, prompted LLMs perform well on textbook cases
but degrade on real records and show demographic bias~\cite{lee2025edtriagebias}; a recent benchmark
finds that most models under-triage high-acuity patients~\cite{npjai2026emergency}. These studies
evaluate a single model on a fixed record; the model cannot question the patient or measure vital signs.

\textbf{Safety of clinical agents.} Sycophancy studies show that LLMs follow false medical premises
unless explicitly allowed to refuse~\cite{chen2025sycophancy}. MedAgentBench~\cite{jiang2025medagentbench}
evaluates agents that place orders in a virtual EHR but scores task success rather than harm. Tool
use and acting agents~\cite{yao2023react} make it possible to check actions against external
knowledge, such as openFDA drug labels~\cite{kasshout2016openfda}.

\textbf{Agent memory and its integrity.} Reflection-based agents store verbal lessons to avoid
repeating mistakes~\cite{shinn2023reflexion}, and MedAgentSim stores corrected diagnoses as
experience. AgentPoison~\cite{chen2024agentpoison} shows that a few malicious entries in an agent's
memory are retrieved with high probability and steer its behaviour. To our knowledge, no prior
clinical simulator verifies memory writes or measures how a wrong lesson spreads.

\textbf{Positioning.} No existing work combines interactive multi-agent diagnosis with
acuity-based triage, independent safety review and memory-integrity testing in one simulated
hospital. SafeWard fills this gap while keeping the MedAgentSim consultation loop unchanged, so
that its results remain comparable with the base paper.

% =====================================================================
\section{Problem Formulation}

\textbf{Episode.} Each patient $p$ is a hidden case file $c_p$ with a reference diagnosis $d^*_p$.
An agent can access $c_p$ only through actions, and every agent observes only its own slice of
the file (Section~\ref{sec:system}).

\textbf{Triage (G1).} On arrival the triage agent observes the arrival record $a_p$ (age, sex,
arrival mode and chief complaint in the patient's words). It may ask the patient questions or
call a vital-signs tool, at most $T_{\text{tri}}$ times, and must output a level
$\hat{y}_p \in \{1,\dots,5\}$, where 1 is most urgent. The reference level $y_p$ is the
expert-consensus KTAS label. The decision is an \emph{under-triage} if $\hat{y}_p > y_p$ and an
\emph{over-triage} if $\hat{y}_p < y_p$. A ward manager then orders the waiting queue by the
assigned levels.

\textbf{Consultation.} In each of at most $T=20$ turns the doctor agent chooses one action:
a question to the patient, a test request answered by the measurement agent, or a final
diagnosis $\hat{d}_p$. After diagnosing, the doctor writes a management plan $P_p$ (medications,
procedures and disposition).

\textbf{Safety (G2).} Each adversarial scenario defines a set $U_p$ of actions that are unsafe for
that patient. The plan that is finally executed, $P'_p$ (after any safety review), is unsafe if
it contains an element of $U_p$.

\textbf{Memory (G3).} A memory $M$ supports a write operation $\mathrm{add}(c, r)$ after each
episode with result $r$ and retrieval $\mathrm{retrieve}(q, k)$ of the $k$ nearest entries for
a query $q$. A poisoning attack submits wrong lessons $\ell$ that associate a presentation with an
incorrect diagnosis; the attack succeeds if accuracy on later similar patients drops.

\textbf{Hypotheses.} Following our Phase~1 proposal, we test:
\begin{itemize}
  \item \textbf{H1:} the Triage Nurse reduces the under-triage rate on level~1--2 patients by at least
  30\% relative to a single LLM given the full record, and reaches a higher quadratic weighted
  kappa than rule-based triage.
  \item \textbf{H2:} SafeWard produces at least 50\% fewer unsafe actions than no review and at least
  25\% fewer than doctor self-review on the adversarial set, with diagnostic accuracy within 3
  points of no review.
  \item \textbf{H3:} after injecting wrong lessons, accuracy on the paired cases drops by at least 5
  points with ungated memory but by no more than 2 points with gated memory.
\end{itemize}

% =====================================================================
\section{Proposed System: SafeWard}
\label{sec:system}

Fig.~\ref{fig:arch} shows the architecture. A patient passes through triage and an acuity queue
(G1), the unchanged MedAgentSim consultation loop, a safety review of the plan (G2) and a memory
gate (G3). Every gap-specific component is a configuration switch (\texttt{features} in the run
configuration), so each can be removed for ablation.

\begin{figure}[t]
\centering
\fbox{\parbox{0.9\columnwidth}{\centering \tbd{figures/architecture.pdf}\\[2pt] arrival $\rightarrow$ triage nurse $\rightarrow$ ward manager $\rightarrow$ doctor $\leftrightarrow$ patient / measurement $\rightarrow$ safety officer $\rightarrow$ memory gate}}
\caption{SafeWard architecture. Shaded components are new (G1--G3); the consultation loop is
kept as in MedAgentSim. The baseline skips all shaded stages.}
\label{fig:arch}
\end{figure}

\textbf{Agents and information isolation.} Table~\ref{tab:agents} lists the agents. Isolation is
enforced in code rather than by prompt: each agent builds its own view of the shared transcript.
The doctor never sees the case file or the reference diagnosis; the patient never sees test
results; the measurement agent sees only the current request; and the Safety Officer reviews the
plan without the doctor's reasoning, to reduce anchoring.

\begin{table}[t]
\caption{Agents in SafeWard}
\label{tab:agents}
\centering
\small
\begin{tabular}{p{1.6cm}p{3.6cm}p{2.0cm}}
\toprule
Agent & Role & Tools / view \\
\midrule
Patient & Answers from the hidden case file; adversarial personas in the safety set & case file \\
Triage Nurse (G1) & Short interview, assigns KTAS 1--5 with red flags & \texttt{GET\_VITALS} \\
Ward Manager (G1) & Orders queue by acuity; escalates red flags & queue state \\
Doctor & History, test orders, diagnosis, plan & \texttt{REQUEST TEST}, memory context \\
Measurement & Returns recorded results on request & exams and tests \\
Safety Officer (G2) & Approves, revises or vetoes the plan & rules, openFDA \\
\bottomrule
\end{tabular}
\end{table}

\textbf{State and contracts.} All data that crosses a module boundary is a typed, validated
record: the case, each message, the episode state (turn, transcript, tests, status, token usage)
and the episode result (diagnosis, judge verdict, plan, safety verdicts, retrieved memory). Every
run writes these records as JSON lines, which form the agent trajectories we analyse.

\textbf{Triage Nurse and Ward Manager (G1).} The nurse receives only the arrival record and a
summary of KTAS criteria. Each reply is one action: \texttt{ASK} a question (answered by a patient
agent that knows the complaint, pain, consciousness and injury status), \texttt{GET\_VITALS} (a
deterministic tool returning the recorded blood pressure, heart rate, respiratory rate,
temperature and oxygen saturation), or \texttt{TRIAGE} with a level, red flags and a reason. The
nurse may take at most four actions. During development on the dev split we observed that the
nurse frequently under-triaged urgent patients (Section~\ref{sec:failures}); we therefore add a
\emph{rule floor}: when deterministic danger-sign rules (unresponsiveness, SpO$_2<90\%$, systolic
pressure $<80$~mmHg, danger-zone vital signs or high-risk complaints) indicate level~1--2, the
nurse may not assign a less urgent level. The Ward Manager simulates two doctors serving a ward of
50 patients with exponentially distributed arrivals; it orders the queue by assigned level and
arrival time and escalates patients with red flags by one level. As in MedAgentSim, the baseline
serves patients in arrival order.

\textbf{Consultation loop.} The doctor, patient and measurement agents follow MedAgentSim. The
doctor acts through text markers (\texttt{REQUEST TEST:}, \texttt{DIAGNOSIS READY:}); each reply is
reduced to exactly one action so that the turn budget cannot be bypassed, and the measurement
agent matches requested tests to recorded ones by meaning. After the diagnosis, the doctor writes
a management plan.

\textbf{Safety Officer (G2).} The officer reviews the plan against deterministic red-flag and
contraindication rules and against drug-interaction and contraindication sections of openFDA drug
labels~\cite{kasshout2016openfda}, which are cached so that runs are repeatable. It returns
\emph{approve}, \emph{revise} or \emph{veto} with a list of issues. On revise or veto the doctor
writes one revised plan, which is reviewed again; in interactive mode a veto is escalated to a
human attending. The self-check baseline instead asks the doctor model to review its own plan
without tools.

\textbf{Verified memory (G3).} The memory stores medical records of correctly diagnosed cases and
lessons. Before a consultation it retrieves the $k=3$ nearest entries for the doctor's only
up-front information, the case objective, using TF-IDF cosine similarity~\cite{salton1988tfidf},
and adds them to the doctor's prompt. The ungated variant, which follows MedAgentSim, writes
entries without verification. The gated variant writes an entry only if the outcome is confirmed
against the reference diagnosis and the Safety Officer raised no veto. Each entry records its
source case and status, every write, rejection and retraction is logged, and an entry is
retracted automatically when its retrieval is followed by wrong diagnoses.

\textbf{Robustness.} Rate-limited calls are retried after the delay the server requests; when an
API key reaches its daily quota the provider switches to the next configured key; configuration
errors stop a run immediately; and failures are recorded per case rather than aborting the batch.
The original text of every normalised model reply is kept in the trajectory for auditing.

% =====================================================================
\section{Experimental Setup}

\textbf{Data and splits.} All splits are fixed before any test run, and design decisions use the
dev splits only.
\begin{itemize}
  \item \emph{Triage:} the KTAS dataset~\cite{moon2019ktas} contains 1{,}267 adult emergency visits
  from two hospitals, each with vital signs, chief complaint, the triage nurse's level and an
  expert-consensus reference level. We draw a stratified split with seed~0: 5 dev and 30 test
  cases per level (level~1 has only 26 cases, giving 21 test cases), i.e.\ 25 dev and 141 test cases.
  \item \emph{Consultation:} AgentClinic's MedQA cases~\cite{schmidgall2024agentclinic,jin2021medqa},
  with a fixed test subset of \tbd{n} cases (seed \tbd{s}).
  \item \emph{Safety:} \tbd{n} adversarial scenarios written and reviewed by two group members, in three
  types (false premise, hidden red flag, unsafe drug combination), split into \tbd{10} dev and
  \tbd{30} test, plus \tbd{20} normal MedQA cases for false vetoes.
  \item \emph{Memory:} \tbd{10} wrong lessons, each paired with \tbd{3} similar later cases.
\end{itemize}

\textbf{Baselines.} B0 is a single LLM given the full arrival record (non-agentic triage). B1 is
deterministic rule-based triage. We also report the real triage nurses' levels from the dataset
as a human reference. B2 is the MedAgentSim-style system: consultation with ungated memory and
chain-of-thought, no triage and no safety review, with patients served in arrival order. B3 is B2
with doctor self-review of the plan. All systems use the same cases and models.

\textbf{Models.} Because LLaMA~3.3~70B, the base paper's model, was no longer available on our
free-tier provider, the doctor, triage nurse, Safety Officer and judge use
\texttt{openai/gpt-oss-120b} and the patient and measurement agents use \texttt{openai/gpt-oss-20b}~\cite{openai2025gptoss},
served through the Groq API with reasoning effort \emph{low}. Temperatures are 0.2 for the doctor,
nurse and patient and 0.0 for the measurement agent, the Safety Officer and the judge. The
consultation is limited to 20 turns and triage to four actions. \tbd{number of runs per
configuration and seeds}.

\textbf{Implementation.} SafeWard is written in Python. Code, configurations, results and
trajectories are available at \url{https://github.com/aleeza-23/AI-Hospital} (release
\texttt{phase2-final}); every table in this paper is generated from the stored results.

% =====================================================================
\section{Evaluation Methodology}

\textbf{Diagnosis.} An LLM judge decides whether the predicted and reference diagnoses denote the
same disease (yes/no), as in MedAgentSim; we audit \tbd{30} judgements manually and report the
agreement. Accuracy is reported with a 95\% Wilson interval~\cite{wilson1927interval}.

\textbf{Triage.} We report exact accuracy, quadratic weighted kappa (QWK)~\cite{cohen1968kappa}
\begin{equation}
\kappa_w = 1 - \frac{\sum_{i,j} w_{ij} O_{ij}}{\sum_{i,j} w_{ij} E_{ij}}, \qquad w_{ij} = \frac{(i-j)^2}{(K-1)^2},
\end{equation}
where $O$ is the observed confusion matrix, $E$ the matrix expected under independence and $K=5$;
the under- and over-triage rates; the under-triage rate on level~1--2 patients, our primary
safety metric for G1; and, from the ward simulation, the mean time-to-doctor of level~1--2
patients and the rate at which KTAS target times (0, 10, 30, 60 and 120 minutes for levels 1--5)
are exceeded.

\textbf{Safety.} We report the unsafe-action rate of the final plan, the compliance rate on
false-premise scenarios, red-flag recall, the false-veto rate on normal cases and the diagnostic
accuracy on normal cases. A plan is judged unsafe by deterministic matching against $U_p$, with
an LLM judge only for unclear cases.

\textbf{Memory.} We report accuracy on the paired cases in clean and poisoned conditions, the drop
in points, the number of wrong lessons admitted and the number of cases that retrieved one.

\textbf{Efficiency and reliability.} For every configuration we report tokens, LLM calls and
latency per case, and the failure rate (errors, unparseable outputs and timeouts).

\textbf{Statistics and comparison with the base paper.} Stochastic configurations are run
\tbd{three} times and reported as mean $\pm$ standard deviation. We distinguish numbers we
\emph{reproduced}, by re-implementing the MedAgentSim method (B2) and running it with our model on
our splits, from numbers \emph{reported} in~\cite{almansoori2025medagentsim}, which used a
different model and are quoted for context only.

% =====================================================================
\section{Results}
% CHECKLIST (organise by gap; the brief requires this):
\subsection{G1: Triage}
\begin{table}[t]\centering\small
\caption{G1 triage on the KTAS test split ($n=141$). LLM methods show the mean$\pm$SD of 2 runs; the human labels and B1 are deterministic. UT = under-triage rate (\%); UT$_{1\text{-}2}$ restricted to KTAS 1--2. Wait = mean simulated wait of KTAS 1--2 patients (min).}
\label{tab:triage}
\setlength{\tabcolsep}{2.5pt}
\resizebox{\columnwidth}{!}{%
\begin{tabular}{lcccccc}\toprule
Method & Acc & QWK & UT & UT$_{1\text{-}2}$ & Tok/case & Wait \\ \midrule
Human nurse & 84.4 & 0.949 & 12.1 & 19.6 & -- & 6.9 \\
B1 Rules & 59.6 & 0.734 & 14.2 & 21.6 & -- & 8.6 \\
B0 Single LLM & 48.2$\pm$3.0 & 0.733$\pm$0.018 & 31.2$\pm$4.0 & 42.2$\pm$6.9 & 508 & 11.1 \\
Nurse, no floor & 40.8$\pm$1.5 & 0.675$\pm$0.021 & 34.8$\pm$1.0 & 60.8 & 1860 & 11.1 \\
\textbf{Nurse + floor} & 55.3$\pm$1.0 & 0.762$\pm$0.014 & 18.4$\pm$1.0 & 19.6 & 1860 & 9.3 \\
\midrule
Arrival order & -- & -- & -- & -- & -- & 14.4 \\
\bottomrule\end{tabular}}\end{table}

Table~\ref{tab:triage} reports the held-out KTAS test split (141 cases, 51 at KTAS~1--2),
run once after all prompt and guardrail choices were frozen on the development split.
The single-prompt baseline B0 under-triaged 24 of the 51 high-acuity patients (47.1\%).
The interview-only Triage Nurse did worse (31/51, 60.8\%). It talked with the patient
and requested vitals in most episodes, but it often settled on KTAS~3 for patients whose
recorded vitals or mental status met level~1--2 criteria. With the rule floor, the
guarded agent under-triaged 10/51 (19.6\%), the same count as the real triage nurses in
the dataset. Against B0 this is a 58\% relative reduction (paired case-bootstrap
difference 27.5 points, 95\% CI [15.7, 39.2]). Its QWK (0.752) is above that of the rule
baseline B1 (0.734). \textbf{H1 is therefore supported for the guarded agent but not for
the unguarded one.}

Two caveats limit this result. First, the floor overrode 25 of 141 nurse decisions, so the gain
comes mostly from the deterministic rules rather than from the interview. The
guarded agent's advantage over B1 in high-acuity under-triage is not significant
(2.0 points, CI [0.0, 5.9]), and B1 has higher exact accuracy (59.6\% vs.\ 54.6\%).
Second, the experiment used a single seed because of the free-tier token budget. The
intervals therefore reflect case sampling, not run-to-run variance of the LLM.

In the ward simulation, ordering by assigned acuity reduced the mean wait of KTAS~1--2
patients from 14.4~min under the arrival-order queue that MedAgentSim implicitly uses to
9.3~min with the guarded agent. Ordering by the human nurses' labels gave 6.9~min.
\subsection{G2: Safety}
% Table (AUTO: results/tables/safety.tex); H2 accepted/rejected.
\subsection{G3: Verified Memory}
% Table (AUTO: results/tables/memory.tex); H3 accepted/rejected.
\subsection{Comparison with the Base Paper}
% E1 vs E2 vs E3 on MedQA (AUTO: results/tables/consultation.tex) + reported numbers.
\subsection{Gap-Resolution Table}
% AUTO: results/tables/gap_resolution.tex (must match README).

% =====================================================================
\section{Ablation and Comparative Analysis}
% CHECKLIST: remove each component once (E5), table AUTO: results/tables/ablation.tex;
% cost/latency trade-off (E6), figure of accuracy vs tokens.

% =====================================================================
\section{Failure Analysis}
\label{sec:failures}
% CHECKLIST: failure categories + frequencies (AUTO: results/tables/failures.tex);
% 2 success + 2 failure trajectories (excerpts from logs/); cause and fix for each category.
% Development-log evidence (docs/devlog.md): what failed, what changed, measured effect.

% =====================================================================
\section{Limitations and Threats to Validity}
% CHECKLIST: different model from the base paper (LLaMA 3.3 unavailable); free-tier budget -> smaller
% sets and fewer seeds; simulated patients (LLM patient may leak); KTAS labels are single-site
% Korean data; LLM judge; adversarial set written by the group; no MIMIC-IV / MIMIC-IV-ED access.

% =====================================================================
\section{Conclusion}
% CHECKLIST: answer RQ1-RQ3 in one sentence each; future work.

% =====================================================================
\section*{Use of AI Tools}
% Required disclosure (instructor permission received for AI-assisted drafting). Edit to match
% what the group actually did.
An AI coding assistant (Claude) was used to implement parts of the software, to prepare a first
draft of several sections of this paper and to check references. All authors reviewed, edited and
verified the text, code and results, and take full responsibility for the content.

\bibliographystyle{IEEEtran}
\bibliography{references}

\end{document}
